\documentclass{article}

% Ćirilica: kodiranje fonta i jezik (srpski, ćirilica)
\usepackage[T2A]{fontenc}
\usepackage[utf8]{inputenc}
\usepackage[serbianc]{babel}

% Veličina strane i margine
\usepackage[a4paper,top=2cm,bottom=2cm,left=3cm,right=3cm,marginparwidth=1.75cm]{geometry}

\usepackage[european]{circuitikz}

% Korisni paketi
\usepackage{amsmath}
\usepackage{graphicx}
\usepackage{float}
\usepackage{booktabs}
\usepackage{tcolorbox}
\tcbuselibrary{breakable}
\usepackage[colorlinks=true, allcolors=blue]{hyperref}

% Цртежи (TikZ)
\usepackage{tikz}
\usetikzlibrary{arrows.meta, calc, decorations.markings, babel}

% TikZ стилови који се понављају на цртежима
\tikzset{
    naboj/.style={circle, draw, thick, minimum size=6mm, inner sep=0pt, font=\small},
    jezgro/.style={circle, draw, thick, fill=gray!15, minimum size=7mm, inner sep=0pt, font=\small},
    elektron/.style={circle, draw, fill=cyan!30, inner sep=0pt, minimum size=5pt},
    supljina/.style={circle, draw, thick, minimum size=5mm, inner sep=0pt, font=\small},
    sredstrelica/.style={postaction={decorate}, decoration={markings, mark=at position 0.55 with {\arrow{Stealth}}}},
    % --- графици наизменичних величина ---
    pomocna/.style={dashed, gray!70, thin},
    mreza/.style={gray!25, very thin},
    kriva/.style={very thick, red!75!black},
    % --- проводник у пресеку: празан, тачка (ка нама), крстић (од нас) ---
    provodnik/.style={circle, draw, thick, fill=white, minimum size=8pt, inner sep=0pt},
    provodnikT/.style={provodnik, path picture={\fill[black] (path picture bounding box.center) circle (1.1pt);}},
    provodnikK/.style={provodnik, path picture={
        \draw (path picture bounding box.south west) -- (path picture bounding box.north east)
              (path picture bounding box.north west) -- (path picture bounding box.south east);}},
    oznaka/.style={fill=white, inner sep=1pt, font=\small},
}

% Оквир за решене примере
\newtcolorbox[auto counter]{primer}{breakable,
    colback=white, colframe=black!55, colbacktitle=black!55, coltitle=white,
    fonttitle=\bfseries\small, title={Пример~\thetcbcounter},
    boxrule=0.5pt, arc=1.5pt, left=6pt, right=6pt, top=3pt, bottom=3pt,
    before skip=8pt, after skip=10pt}

% Оквир за кључне формуле
\newtcolorbox{kljucno}{colback=blue!4, colframe=blue!45!black,
    boxrule=0.5pt, arc=1.5pt, left=8pt, right=8pt, top=4pt, bottom=4pt}

\title{ОЕТ 1 - Вежбање}
\author{Саобраћајно техничка школа, Београд, Земун}

\begin{document}


\section{Основне формуле}

\subsection{Електрична отпорност}

Отпорност проводника зависи од његове дужине $l$, површине попречног пресека $S$
и материјала од ког је направљен (специфична отпорност $\rho$):
\begin{kljucno}
\[
R = \rho \, \frac{l}{S} \qquad [\Omega]
\]
\end{kljucno}

Проводност је реципрочна вредност отпорности:
\[
G = \frac{1}{R} \qquad [\text{S}]
\]

Отпорност се мења са температуром ($\alpha$ је температурни коефицијент):
\[
R = R_0 \left(1 + \alpha \, \Delta t\right)
\]

\textbf{Редна веза} -- отпорности се сабирају:
\[
R_e = R_1 + R_2 + \dots + R_n
\]

\textbf{Паралелна веза} -- сабирају се реципрочне вредности:
\[
\frac{1}{R_e} = \frac{1}{R_1} + \frac{1}{R_2} + \dots + \frac{1}{R_n},
\qquad \text{за два отпорника:} \quad R_e = \frac{R_1 R_2}{R_1 + R_2}
\]

\subsection{Омов закон и снага}

Струја кроз отпорник сразмерна је напону, а обрнуто сразмерна отпорности:
\begin{kljucno}
\[
I = \frac{U}{R} \qquad \Longrightarrow \qquad U = I R, \qquad R = \frac{U}{I}
\]
\end{kljucno}

Снага на отпорнику:
\[
P = U I = I^2 R = \frac{U^2}{R} \qquad [\text{W}]
\]

\subsection{Капацитивност}

Капацитивност показује колико наелектрисања $Q$ кондензатор прими при напону $U$:
\begin{kljucno}
\[
C = \frac{Q}{U} \qquad [\text{F}]
\]
\end{kljucno}

Капацитивност плочастог кондензатора зависи од површине плоча $S$, растојања
између њих $d$ и диелектрика између плоча:
\[
C = \varepsilon_0 \, \varepsilon_r \, \frac{S}{d},
\qquad \varepsilon_0 = 8{,}85 \cdot 10^{-12}\,\frac{\text{F}}{\text{m}}
\]

\textbf{Паралелна веза} -- капацитивности се сабирају:
\[
C_e = C_1 + C_2 + \dots + C_n
\]

\textbf{Редна веза} -- сабирају се реципрочне вредности:
\[
\frac{1}{C_e} = \frac{1}{C_1} + \frac{1}{C_2} + \dots + \frac{1}{C_n},
\qquad \text{за два кондензатора:} \quad C_e = \frac{C_1 C_2}{C_1 + C_2}
\]

Енергија напуњеног кондензатора:
\[
W = \frac{C U^2}{2} = \frac{Q U}{2} \qquad [\text{J}]
\]

%==================================================================
% БЕЛЕШКЕ ИЗ СВЕСКЕ (редом по сликама IMG_6501 -- IMG_6508)
%==================================================================

\end{document}